\chapter{Why Logging Matters}
Logging is how a system narrates events to the people who operate it. A good log message explains the event name, the outcome, and the key fields and IDs. It should do this fast, without drama. The goal is simple: make investigation cheaper.

\section{What a Good Log Gives You}
A strong log message helps you answer three questions:
\begin{itemize}
\item What event happened?
\item What was the outcome and level?
\item Which IDs or fields let me trace it?
\end{itemize}

When those answers are present, the message is actionable. When they are missing, every incident becomes a hunt.

\section{The Core Vocabulary}
Use consistent terms across the book and across services. This keeps team conversations short and precise.

\begin{table}[h]
\centering
\begin{tabularx}{\textwidth}{lX}
\toprule
Term & Meaning \\
\midrule
Event name & A short verb phrase that describes the action, for example \texttt{job.completed}. \\
Level & The severity filter: debug, info, warn, error, or fatal. \\
Outcome & The result of the event: success, failure, or retry. \\
Fields & Structured data attached to the event. \\
IDs & Stable identifiers such as request id, correlation id, user id, or job id. \\
\bottomrule
\end{tabularx}
\end{table}

\section{Checklist: The Minimum Message}
\begin{itemize}
\item Event name is explicit.
\item Outcome is stated.
\item Level matches the outcome.
\item IDs are present and stable.
\item Fields cover the minimum context.
\end{itemize}

